\input{preamble/preamble.tex}

\geometry{margin=0.85in}
\AtBeginDocument{\small}

\newcommand{\ExamNameField}{\noindent\textbf{Name:}\ \rule{0.7\linewidth}{0.4pt}\par\medskip}

\begin{document}
	\ExamSection{C9: Confidence interval for the mean when variance is unknown}
	\begin{ExamProblems}
		\item
		\subsection*{Problem description}
		A microfinance institution studies daily loan disbursement amounts (hundreds of USD) in one branch.
		A random sample of $n=10$ days gives the raw data:
		\[
		42,\ 38,\ 45,\ 40,\ 41,\ 39,\ 44,\ 43,\ 37,\ 41
		\]
		Construct both a 95\% t-confidence interval and a 95\% z-confidence interval for the population mean, then compare them.

		\subsection*{C9}
		Step 1: Identify parameter and compute sample statistics
		\[
		\mu=\text{population mean daily loan disbursement},\quad n=10
		\]
		\[
		\bar{x}=\frac{\sum x_i}{n}=\frac{410}{10}=41.00
		\]
		\[
		\sum (x_i-\bar{x})^2=60,\quad s=\sqrt{\frac{\sum (x_i-\bar{x})^2}{n-1}}=\sqrt{\frac{60}{9}}=2.582
		\]
		\[
		df=n-1=9
		\]
		Step 2: Construct t-interval
		\[
		SE=\frac{s}{\sqrt{n}}=\frac{2.582}{\sqrt{10}}=0.816,\qquad t^*=t_{0.025,9}\approx 2.262
		\]
		\[
		\mu\in \bar{x}\pm t^*SE=41.00\pm 2.262(0.816)=41.00\pm 1.85=[39.15,\,42.85]
		\]
		Step 3: Construct z-interval
		\[
		z^*=z_{0.025}=1.960,\qquad \mu\in 41.00\pm 1.960(0.816)=41.00\pm 1.60=[39.40,\,42.60]
		\]
		Step 4: Numerical comparison
		\[
		t^*=2.262>1.960=z^*,\quad E_t=1.85>1.60=E_z
		\]
		\[
		\text{Width}_t=2(1.85)=3.70,\qquad \text{Width}_z=2(1.60)=3.20
		\]
		\[
		\text{Absolute difference}=\text{Width}_t-\text{Width}_z=3.70-3.20=0.50
		\]
		\[
		\text{Relative difference (base }z\text{-interval)}=\frac{3.70-3.20}{3.20}\times 100=15.63\%
		\]
		The t-interval is 15.63\% wider than the z-interval, using the z-interval width as the base reference.
		Step 5: Conceptual explanation of convergence
		With only $df=9$, the t critical value is noticeably larger than z, so the t-interval is visibly wider.
		As $df$ increases, $t(df)\to N(0,1)$ and both methods get closer.
		Step 6: Interpretation in economic/financial context
		The branch's average daily disbursement is estimated near 41 hundred USD.
		Using the exact method for unknown variance, the 95\% t-confidence interval is $[39.15,42.85]$ hundred USD.

		\newpage
		\item
		\subsection*{Problem description}
		A brokerage firm analyzes mean transaction fees (USD) charged to retail investors.
		A random sample of $n=18$ transactions gives the raw data:
		\[
		1.8,\ 2.1,\ 1.9,\ 2.4,\ 2.0,\ 1.7,\ 2.2,\ 2.3,\ 1.8,\ 2.5,\ 2.1,\ 1.9,\ 2.0,\ 2.2,\ 1.6,\ 2.4,\ 2.3,\ 1.9
		\]
		Construct both a 90\% t-confidence interval and a 90\% z-confidence interval for the population mean, then compare them.

		\subsection*{C9}
		Step 1: Identify parameter and compute sample statistics
		\[
		\mu=\text{population mean transaction fee},\quad n=18
		\]
		\[
		\bar{x}=\frac{\sum x_i}{n}=\frac{37.1}{18}=2.061
		\]
		\[
		\sum (x_i-\bar{x})^2=1.143,\quad s=\sqrt{\frac{1.143}{17}}=0.259
		\]
		\[
		df=17
		\]
		Step 2: Construct t-interval
		\[
		SE=\frac{0.259}{\sqrt{18}}=0.0611,\qquad t^*=t_{0.05,17}\approx 1.740
		\]
		\[
		\mu\in 2.061\pm 1.740(0.0611)=2.061\pm 0.106=[1.955,\,2.167]
		\]
		Step 3: Construct z-interval
		\[
		z^*=z_{0.05}=1.645,\qquad \mu\in 2.061\pm 1.645(0.0611)=2.061\pm 0.101=[1.960,\,2.162]
		\]
		Step 4: Numerical comparison
		\[
		t^*=1.740>1.645=z^*,\quad E_t=0.106>E_z=0.101
		\]
		\[
		\text{Width}_t=0.212,\qquad \text{Width}_z=0.202
		\]
		\[
		\text{Absolute difference}=0.212-0.202=0.010
		\]
		\[
		\text{Relative difference (base }z\text{-interval)}=\frac{0.212-0.202}{0.202}\times 100=4.95\%
		\]
		The t-interval is 4.95\% wider than the z-interval, using the z-interval width as the base reference.
		Step 5: Conceptual explanation of convergence
		At $df=17$, t is still distinctly heavier-tailed than z, so the exact t-interval remains clearly wider.
		Increasing sample size would reduce this gap.
		Step 6: Interpretation in economic/financial context
		The firm's mean transaction fee is around 2.0 USD.
		Using the exact t method, the 90\% confidence interval is $[1.955,2.167]$ USD.

		\newpage
		\item
		\subsection*{Problem description}
		A supermarket chain studies mean customer basket value (USD) during weekday mornings.
		A random sample of $n=30$ receipts gives the raw data:
		\[
		115,\ 122,\ 118,\ 130,\ 126,\ 119,\ 121,\ 124,\ 128,\ 117,
		123,\ 120,\ 125,\ 129,\ 116,\ 127,\ 122,\]
		\[	 118,\ 124,\ 121,
		126,\ 119,\ 123,\ 128,\ 120,\ 125,\ 117,\ 124,\ 122,\ 126
		\]
		Construct both a 99\% t-confidence interval and a 99\% z-confidence interval for the population mean, then compare them.

		\subsection*{C9}
		Step 1: Identify parameter and compute sample statistics
		\[
		\mu=\text{population mean basket value},\quad n=30
		\]
		\[
		\bar{x}=\frac{3675}{30}=122.50
		\]
		\[
		\sum (x_i-\bar{x})^2=477.5,\quad s=\sqrt{\frac{477.5}{29}}=4.058
		\]
		\[
		df=29
		\]
		Step 2: Construct t-interval
		\[
		SE=\frac{4.058}{\sqrt{30}}=0.741,\qquad t^*=t_{0.005,29}\approx 2.756
		\]
		\[
		\mu\in 122.50\pm 2.756(0.741)=122.50\pm 2.04=[120.46,\,124.54]
		\]
		Step 3: Construct z-interval
		\[
		z^*=z_{0.005}=2.576,\qquad \mu\in 122.50\pm 2.576(0.741)=122.50\pm 1.91=[120.59,\,124.41]
		\]
		Step 4: Numerical comparison
		\[
		t^*=2.756>2.576=z^*,\quad E_t=2.04>1.91=E_z
		\]
		\[
		\text{Width}_t=4.08,\qquad \text{Width}_z=3.82
		\]
		\[
		\text{Absolute difference}=4.08-3.82=0.26
		\]
		\[
		\text{Relative difference (base }z\text{-interval)}=\frac{4.08-3.82}{3.82}\times 100=6.81\%
		\]
		The t-interval is 6.81\% wider than the z-interval, using the z-interval width as the base reference.
		Step 5: Conceptual explanation of convergence
		At $n=30$, the sample is already moderate, so t and z are closer than in very small samples.
		Still, the exact t method remains preferred because variance is unknown.
		Step 6: Interpretation in economic/financial context
		The chain's weekday-morning average basket value is estimated around 121--125 USD.
		Using the exact method, the 99\% t-confidence interval is $[120.46,124.54]$ USD.

		\newpage
		\item
		\subsection*{Problem description}
		An e-commerce platform studies average delivery time (hours) for same-city orders.
		For $n=60$ orders, data are grouped as follows:
		\[
		\begin{array}{c|c|c}
		\text{Class (hours)} & \text{Midpoint }m_i & f_i\\\hline
		20\text{--}24 & 22 & 6\\
		25\text{--}29 & 27 & 9\\
		30\text{--}34 & 32 & 14\\
		35\text{--}39 & 37 & 13\\
		40\text{--}44 & 42 & 10\\
		45\text{--}49 & 47 & 8
		\end{array}
		\]
		Construct both a 95\% t-confidence interval and a 95\% z-confidence interval for the population mean, then compare them.

		\subsection*{C9}
		Step 1: Identify parameter and compute grouped statistics
		\[
		\mu=\text{population mean delivery time},\quad n=\sum f_i=60
		\]
		\[
		\sum f_im_i=2100,\quad \sum f_im_i^2=77910
		\]
		\[
		\bar{x}\approx\frac{\sum f_im_i}{n}=\frac{2100}{60}=35.00
		\]
		\[
		s\approx\sqrt{\frac{\sum f_im_i^2-n\bar{x}^2}{n-1}}=\sqrt{\frac{77910-60(35)^2}{59}}=7.602
		\]
		\[
		df=59
		\]
		Step 2: Construct t-interval
		\[
		SE=\frac{7.602}{\sqrt{60}}=0.981,\qquad t^*=t_{0.025,59}\approx 2.000
		\]
		\[
		\mu\in 35.00\pm 2.000(0.981)=35.00\pm 1.96=[33.04,\,36.96]
		\]
		Step 3: Construct z-interval
		\[
		z^*=1.960,\qquad \mu\in 35.00\pm 1.960(0.981)=35.00\pm 1.92=[33.08,\,36.92]
		\]
		Step 4: Numerical comparison
		\[
		t^*=2.000>1.960=z^*,\quad E_t=1.96>1.92=E_z
		\]
		\[
		\text{Width}_t=3.92,\qquad \text{Width}_z=3.84
		\]
		\[
		\text{Absolute difference}=3.92-3.84=0.08
		\]
		\[
		\text{Relative difference (base }z\text{-interval)}=\frac{3.92-3.84}{3.84}\times 100=2.08\%
		\]
		The t-interval is 2.08\% wider than the z-interval, using the z-interval width as the base reference.
		Step 5: Conceptual explanation of convergence
		With $df=59$, t and z are already very close, so interval widths are similar.
		This illustrates convergence of t to the normal distribution.
		Step 6: Interpretation in economic/financial context
		Average same-city delivery time is estimated near 35 hours.
		Using the exact t method, the 95\% confidence interval is $[33.04,36.96]$ hours.

		\newpage
		\item
		\subsection*{Problem description}
		A central warehouse estimates mean daily inventory holding cost (thousand USD).
		For $n=120$ days, grouped data are:
		\[
		\begin{array}{c|c|c}
		\text{Class (thousand USD)} & m_i & f_i\\\hline
		50\text{--}55 & 52.5 & 12\\
		55\text{--}60 & 57.5 & 18\\
		60\text{--}65 & 62.5 & 28\\
		65\text{--}70 & 67.5 & 24\\
		70\text{--}75 & 72.5 & 22\\
		75\text{--}80 & 77.5 & 16
		\end{array}
		\]
		Construct both a 90\% t-confidence interval and a 90\% z-confidence interval for the population mean, then compare them.

		\subsection*{C9}
		Step 1: Identify parameter and compute grouped statistics
		\[
		\mu=\text{population mean daily inventory holding cost},\quad n=120
		\]
		\[
		\sum f_im_i=7870,\quad \sum f_im_i^2=524975
		\]
		\[
		\bar{x}\approx\frac{7870}{120}=65.583
		\]
		\[
		s\approx\sqrt{\frac{524975-120(65.583)^2}{119}}=7.620
		\]
		\[
		df=119
		\]
		Step 2: Construct t-interval
		\[
		SE=\frac{7.620}{\sqrt{120}}=0.696,\qquad t^*=t_{0.05,119}\approx 1.658
		\]
		\[
		\mu\in 65.583\pm 1.658(0.696)=65.583\pm 1.15=[64.43,\,66.74]
		\]
		Step 3: Construct z-interval
		\[
		z^*=1.645,\qquad \mu\in 65.583\pm 1.645(0.696)=65.583\pm 1.14=[64.44,\,66.72]
		\]
		Step 4: Numerical comparison
		\[
		t^*=1.658>1.645=z^*,\quad E_t=1.15>E_z=1.14
		\]
		\[
		\text{Width}_t=2.30,\qquad \text{Width}_z=2.28
		\]
		\[
		\text{Absolute difference}=2.30-2.28=0.02
		\]
		\[
		\text{Relative difference (base }z\text{-interval)}=\frac{2.30-2.28}{2.28}\times 100=0.88\%
		\]
		The t-interval is 0.88\% wider than the z-interval, using the z-interval width as the base reference.
		Step 5: Conceptual explanation of convergence
		At $df=119$, the critical values are almost equal, so the two intervals nearly coincide.
		This is the expected large-sample behavior.
		Step 6: Interpretation in economic/financial context
		The warehouse's mean daily holding cost is estimated around 65 to 67 thousand USD.
		Using the exact t method, the 90\% confidence interval is $[64.43,66.74]$ thousand USD.

		\newpage
		\item
		\subsection*{Problem description}
		A regional tax office studies the mean monthly VAT payment (thousand USD) of small firms.
		For $n=250$ firms, grouped data are:
		\[
		\begin{array}{c|c|c}
		\text{Class (thousand USD)} & m_i & f_i\\\hline
		80\text{--}90 & 85 & 18\\
		90\text{--}100 & 95 & 32\\
		100\text{--}110 & 105 & 46\\
		110\text{--}120 & 115 & 58\\
		120\text{--}130 & 125 & 44\\
		130\text{--}140 & 135 & 30\\
		140\text{--}150 & 145 & 22
		\end{array}
		\]
		Construct both a 95\% t-confidence interval and a 95\% z-confidence interval for the population mean, then compare them.

		\subsection*{C9}
		Step 1: Identify parameter and compute grouped statistics
		\[
		\mu=\text{population mean monthly VAT payment},\quad n=250
		\]
		\[
		\sum f_im_i=28810,\quad \sum f_im_i^2=3403450
		\]
		\[
		\bar{x}\approx\frac{28810}{250}=115.24
		\]
		\[
		s\approx\sqrt{\frac{3403450-250(115.24)^2}{249}}=16.741
		\]
		\[
		df=249
		\]
		Step 2: Construct t-interval
		\[
		SE=\frac{16.741}{\sqrt{250}}=1.059,\qquad t^*=t_{0.025,249}\approx 1.969
		\]
		\[
		\mu\in 115.24\pm 1.969(1.059)=115.24\pm 2.08=[113.16,\,117.32]
		\]
		Step 3: Construct z-interval
		\[
		z^*=1.960,\qquad \mu\in 115.24\pm 1.960(1.059)=115.24\pm 2.08=[113.16,\,117.32]
		\]
		Step 4: Numerical comparison
		\[
		t^*=1.969>1.960=z^*,\quad E_t=2.08\gtrsim 2.08=E_z
		\]
		\[
		\text{Width}_t=4.16,\qquad \text{Width}_z=4.16
		\]
		\[
		\text{Absolute difference}=4.16-4.16=0.00\text{ (more precisely }0.02\text{)}
		\]
		\[
		\text{Relative difference (base }z\text{-interval)}\approx\frac{0.02}{4.16}\times 100=0.48\%
		\]
		The t-interval is about 0.48\% wider than the z-interval, using the z-interval width as the base reference.
		Step 5: Conceptual explanation of convergence
		For $df=249$, the t and z critical values are extremely close.
		Hence, practical differences in interval width are very small.
		Step 6: Interpretation in economic/financial context
		The mean monthly VAT payment is estimated near 115 thousand USD.
		Using the exact t method, the 95\% confidence interval is $[113.16,117.32]$ thousand USD.

		\newpage
		\item
		\subsection*{Problem description}
		A digital bank estimates the mean monthly balance (hundreds of USD) of active app users.
		For $n=500$ users, grouped data are:
		\[
		\begin{array}{c|c|c}
		\text{Class (hundreds of USD)} & m_i & f_i\\\hline
		60\text{--}65 & 62.5 & 30\\
		65\text{--}70 & 67.5 & 52\\
		70\text{--}75 & 72.5 & 84\\
		75\text{--}80 & 77.5 & 108\\
		80\text{--}85 & 82.5 & 96\\
		85\text{--}90 & 87.5 & 68\\
		90\text{--}95 & 92.5 & 40\\
		95\text{--}100 & 97.5 & 22
		\end{array}
		\]
		Construct both a 99\% t-confidence interval and a 99\% z-confidence interval for the population mean, then compare them.

		\subsection*{C9}
		Step 1: Identify parameter and compute grouped statistics
		\[
		\mu=\text{population mean monthly balance},\quad n=500
		\]
		\[
		\sum f_im_i=39560,\quad \sum f_im_i^2=3174700
		\]
		\[
		\bar{x}\approx\frac{39560}{500}=79.12
		\]
		\[
		s\approx\sqrt{\frac{3174700-500(79.12)^2}{499}}=8.924
		\]
		\[
		df=499
		\]
		Step 2: Construct t-interval
		\[
		SE=\frac{8.924}{\sqrt{500}}=0.399,\qquad t^*=t_{0.005,499}\approx 2.586
		\]
		\[
		\mu\in 79.12\pm 2.586(0.399)=79.12\pm 1.03=[78.09,\,80.15]
		\]
		Step 3: Construct z-interval
		\[
		z^*=2.576,\qquad \mu\in 79.12\pm 2.576(0.399)=79.12\pm 1.03=[78.09,\,80.15]
		\]
		Step 4: Numerical comparison
		\[
		t^*=2.586>2.576=z^*,\quad E_t=1.03\gtrsim 1.03=E_z
		\]
		\[
		\text{Width}_t=2.06,\qquad \text{Width}_z=2.06
		\]
		\[
		\text{Absolute difference}=2.06-2.06=0.00\text{ (more precisely }0.01\text{)}
		\]
		\[
		\text{Relative difference (base }z\text{-interval)}\approx\frac{0.01}{2.06}\times 100=0.49\%
		\]
		The t-interval is about 0.49\% wider than the z-interval, using the z-interval width as the base reference.
		Step 5: Conceptual explanation of convergence
		With hundreds of observations, t and z intervals are nearly indistinguishable.
		This reflects strong convergence of t to normality.
		Step 6: Interpretation in economic/financial context
		The bank's average monthly balance is estimated around 79 to 80 hundred USD.
		Using the exact t method, the 99\% confidence interval is $[78.09,80.15]$ hundred USD.

		\newpage
		\item
		\subsection*{Problem description}
		A national insurer estimates mean annual claim cost (USD) for small commercial policies.
		For $n=1000$ policies, grouped data are:
		\[
		\begin{array}{c|c|c}
		\text{Class (USD)} & m_i & f_i\\\hline
		125\text{--}135 & 130 & 70\\
		135\text{--}145 & 140 & 130\\
		145\text{--}155 & 150 & 220\\
		155\text{--}165 & 160 & 250\\
		165\text{--}175 & 170 & 180\\
		175\text{--}185 & 180 & 100\\
		185\text{--}195 & 190 & 50
		\end{array}
		\]
		Construct both a 95\% t-confidence interval and a 95\% z-confidence interval for the population mean, then compare them.

		\subsection*{C9}
		Step 1: Identify parameter and compute grouped statistics
		\[
		\mu=\text{population mean annual claim cost},\quad n=1000
		\]
		\[
		\sum f_im_i=158400,\quad \sum f_im_i^2=25396000
		\]
		\[
		\bar{x}\approx\frac{158400}{1000}=158.40
		\]
		\[
		s\approx\sqrt{\frac{25396000-1000(158.4)^2}{999}}=15.417
		\]
		\[
		df=999
		\]
		Step 2: Construct t-interval
		\[
		SE=\frac{15.417}{\sqrt{1000}}=0.488,\qquad t^*=t_{0.025,999}\approx 1.962
		\]
		\[
		\mu\in 158.40\pm 1.962(0.488)=158.40\pm 0.96=[157.44,\,159.36]
		\]
		Step 3: Construct z-interval
		\[
		z^*=1.960,\qquad \mu\in 158.40\pm 1.960(0.488)=158.40\pm 0.96=[157.44,\,159.36]
		\]
		Step 4: Numerical comparison
		\[
		t^*=1.962>1.960=z^*,\quad E_t\gtrsim E_z
		\]
		\[
		\text{Width}_t=1.92,\qquad \text{Width}_z=1.92
		\]
		\[
		\text{Absolute difference}\approx 0.00\text{ (more precisely }0.00\text{ to two decimals)}
		\]
		\[
		\text{Relative difference (base }z\text{-interval)}\approx 0.10\%
		\]
		The t-interval is approximately 0.10\% wider than the z-interval, using the z-interval width as the base reference.
		Step 5: Conceptual explanation of convergence
		For $df=999$, convergence is very strong, so t and z decisions are practically identical.
		Step 6: Interpretation in economic/financial context
		Mean annual claim cost is estimated around 158 USD.
		Using the exact t method, the 95\% confidence interval is $[157.44,159.36]$ USD.

		\newpage
		\item
		\subsection*{Problem description}
		A payment network analyzes mean monthly card-processing revenue (thousand USD) per merchant.
		For $n=2500$ merchants, grouped data are:
		\[
		\begin{array}{c|c|c}
		\text{Class (thousand USD)} & m_i & f_i\\\hline
		200\text{--}220 & 210 & 180\\
		220\text{--}240 & 230 & 360\\
		240\text{--}260 & 250 & 620\\
		260\text{--}280 & 270 & 700\\
		280\text{--}300 & 290 & 380\\
		300\text{--}320 & 310 & 180\\
		320\text{--}340 & 330 & 80
		\end{array}
		\]
		Construct both a 90\% t-confidence interval and a 90\% z-confidence interval for the population mean, then compare them.

		\subsection*{C9}
		Step 1: Identify parameter and compute grouped statistics
		\[
		\mu=\text{population mean monthly card-processing revenue},\quad n=2500
		\]
		\[
		\sum f_im_i=657000,\quad \sum f_im_i^2=174410000
		\]
		\[
		\bar{x}\approx\frac{657000}{2500}=262.80
		\]
		\[
		s\approx\sqrt{\frac{174410000-2500(262.8)^2}{2499}}=28.784
		\]
		\[
		df=2499
		\]
		Step 2: Construct t-interval
		\[
		SE=\frac{28.784}{\sqrt{2500}}=0.576,\qquad t^*=t_{0.05,2499}\approx 1.646
		\]
		\[
		\mu\in 262.80\pm 1.646(0.576)=262.80\pm 0.95=[261.85,\,263.75]
		\]
		Step 3: Construct z-interval
		\[
		z^*=1.645,\qquad \mu\in 262.80\pm 1.645(0.576)=262.80\pm 0.95=[261.85,\,263.75]
		\]
		Step 4: Numerical comparison
		\[
		t^*=1.646>1.645=z^*,\quad E_t\gtrsim E_z
		\]
		\[
		\text{Width}_t=1.90,\qquad \text{Width}_z=1.90
		\]
		\[
		\text{Absolute difference}\approx 0.00\text{ (more precisely }0.00\text{ to two decimals)}
		\]
		\[
		\text{Relative difference (base }z\text{-interval)}\approx 0.05\%
		\]
		The t-interval is approximately 0.05\% wider than the z-interval, using the z-interval width as the base reference.
		Step 5: Conceptual explanation of convergence
		At this scale, the sampling distribution behavior makes the two intervals nearly the same.
		The t method remains the exact method under unknown variance.
		Step 6: Interpretation in economic/financial context
		Mean monthly merchant revenue is estimated near 263 thousand USD.
		Using the exact t method, the 90\% confidence interval is $[261.85,263.75]$ thousand USD.

		\newpage
		\item
		\subsection*{Problem description}
		A logistics federation studies mean quarterly freight billing (thousand USD) across firms.
		For $n=5000$ firms, grouped data are:
		\[
		\begin{array}{c|c|c}
		\text{Class (thousand USD)} & m_i & f_i\\\hline
		400\text{--}440 & 420 & 320\\
		440\text{--}480 & 460 & 680\\
		480\text{--}520 & 500 & 1300\\
		520\text{--}560 & 540 & 1400\\
		560\text{--}600 & 580 & 820\\
		600\text{--}640 & 620 & 360\\
		640\text{--}680 & 660 & 120
		\end{array}
		\]
		Construct both a 99\% t-confidence interval and a 99\% z-confidence interval for the population mean, then compare them.

		\subsection*{C9}
		Step 1: Identify parameter and compute grouped statistics
		\[
		\mu=\text{population mean quarterly freight billing},\quad n=5000
		\]
		\[
		\sum f_im_i=2631200,\quad \sum f_im_i^2=1401024000
		\]
		\[
		\bar{x}\approx\frac{2631200}{5000}=526.24
		\]
		\[
		s\approx\sqrt{\frac{1401024000-5000(526.24)^2}{4999}}=55.570
		\]
		\[
		df=4999
		\]
		Step 2: Construct t-interval
		\[
		SE=\frac{55.570}{\sqrt{5000}}=0.786,\qquad t^*=t_{0.005,4999}\approx 2.577
		\]
		\[
		\mu\in 526.24\pm 2.577(0.786)=526.24\pm 2.03=[524.21,\,528.27]
		\]
		Step 3: Construct z-interval
		\[
		z^*=2.576,\qquad \mu\in 526.24\pm 2.576(0.786)=526.24\pm 2.02=[524.22,\,528.26]
		\]
		Step 4: Numerical comparison
		\[
		t^*=2.577>2.576=z^*,\quad E_t\gtrsim E_z
		\]
		\[
		\text{Width}_t=4.06,\qquad \text{Width}_z=4.04
		\]
		\[
		\text{Absolute difference}=4.06-4.04=0.02
		\]
		\[
		\text{Relative difference (base }z\text{-interval)}=\frac{4.06-4.04}{4.04}\times 100=0.50\%
		\]
		The t-interval is 0.50\% wider than the z-interval, using the z-interval width as the base reference.
		Step 5: Conceptual explanation of convergence
		With very large degrees of freedom, t and z are almost identical in practice.
		Hence, the numerical gap is minimal.
		Step 6: Interpretation in economic/financial context
		The federation's average quarterly freight billing is estimated around 526 thousand USD.
		Using the exact t method, the 99\% confidence interval is $[524.21,528.27]$ thousand USD.
	\end{ExamProblems}
\end{document}
